\documentclass[12pt]{article}
\usepackage{amsmath, amssymb, amsthm}
\usepackage{geometry}
\geometry{margin=1in}

\newtheorem{theorem}{Theorem}
\newtheorem{lemma}[theorem]{Lemma}
\newtheorem{corollary}[theorem]{Corollary}
\newtheorem{claim}[theorem]{Claim}
\theoremstyle{definition}
\newtheorem{definition}[theorem]{Definition}
\theoremstyle{remark}
\newtheorem{remark}[theorem]{Remark}

\title{Existence of $\varepsilon$-Light Subsets in Graphs\\[4pt]
\large A Rigorous Proof via the Barrier Function Method}
\author{}
\date{}

\begin{document}
\maketitle

% ================================================================
\section{Setup}

Let $G=(V,E,w)$ be a weighted undirected graph, $n=|V|$.
For $S\subseteq V$, let $G_S$ keep all vertices but only edges with both
endpoints in $S$; write $L,L_S$ for the respective Laplacians.

\begin{definition}
$S$ is \emph{$\varepsilon$-light} if $L_S\preceq\varepsilon L$.
\end{definition}

\noindent Equivalently (quadratic-form reformulation):
\[
  \forall x\in\mathbb{R}^n:\quad
  \sum_{(u,v)\in E(S,S)} w_{uv}(x_u-x_v)^2
  \;\le\;
  \varepsilon\sum_{(u,v)\in E}w_{uv}(x_u-x_v)^2.
\]

\begin{theorem}[Main]\label{thm:main}
There exists a universal constant $c>0$ such that for every weighted graph
$G=(V,E,w)$ and every $\varepsilon\in(0,1]$, there is an $\varepsilon$-light
set $S\subseteq V$ with $|S|\ge c\,\varepsilon n$.
In particular $c=1/42$ is achievable.
\end{theorem}

The proof splits on the \emph{average (weighted) degree}
$d_{\mathrm{avg}}=\frac{2}{n}\sum_{(u,v)\in E}w_{uv}$.

% ================================================================
\section{Sparse Case}

\begin{lemma}[Greedy independent set]\label{lem:ind}
Every graph with $n$ vertices and average degree $d$ has an independent set
of size $\ge n/(d+1)$.
\end{lemma}

\begin{proof}
Process vertices in any order; include vertex $v$ if none of its
already-included neighbours is in $I$.  Every vertex not in $I$ is adjacent
to at least one vertex in $I$, so $n\le|I|+|I|\cdot d_{\mathrm{avg}}=|I|(d+1)$.
\end{proof}

\begin{corollary}\label{cor:sparse}
If $d_{\mathrm{avg}}\le1/(2\varepsilon)$, then there is an $\varepsilon$-light
independent set $I$ with $|I|\ge\varepsilon n/2$.
\end{corollary}

\begin{proof}
By Lemma~\ref{lem:ind}, $|I|\ge n/(d_{\mathrm{avg}}+1)\ge
2\varepsilon n/(1+2\varepsilon)\ge\varepsilon n/2$
(using $1+2\varepsilon\le4$ for $\varepsilon\le1$).
Since $I$ is independent, $L_I=0\preceq\varepsilon L$.
\end{proof}

% ================================================================
\section{Normalised Laplacian and Leverage Scores}

For the dense case we work in the image of $L$.  Let $L^\dagger$ be the
Moore--Penrose pseudoinverse of $L$, and set
$L^{\dagger/2}=(L^\dagger)^{1/2}$.  Define
\[
  \widetilde{L}_S \;=\; L^{\dagger/2}L_S L^{\dagger/2}.
\]
Then $L_S\preceq\varepsilon L$ is equivalent to
$\widetilde{L}_S\preceq\varepsilon I_{\mathrm{im}(L)}$, i.e.,
$\lambda_{\max}(\widetilde{L}_S)\le\varepsilon$.

\noindent The \emph{leverage score} of edge $(s,t)$ is
\[
  \ell(s,t) \;=\; w_{st}(\mathbf{e}_s-\mathbf{e}_t)^\top
               L^\dagger(\mathbf{e}_s-\mathbf{e}_t)
            \;=\; \mathrm{Tr}\!\left(L^{\dagger/2}L_{st}L^{\dagger/2}\right),
\]
with $L_{st}=w_{st}(\mathbf{e}_s-\mathbf{e}_t)(\mathbf{e}_s-\mathbf{e}_t)^\top$.
Key properties: $0\le\ell(s,t)\le1$ and
$\sum_{(s,t)\in E}\ell(s,t)=\mathrm{rank}(L)\le n-1$.

\noindent For a vertex $v$ define $\ell(v)=\sum_{u:(u,v)\in E}\ell(u,v)$
(leverage of all edges incident to $v$) and for sets $A,B$:
\[
  \ell(A,B)=\sum_{a\in A,b\in B,(a,b)\in E}\ell(a,b),\qquad
  \ell(A)=\ell(A,V\setminus A).
\]

% ================================================================
\section{Dense Case: Barrier Function Construction}

\subsection{Parameters}

Fix
\[
  \delta=\frac{21}{n},\quad
  \phi=\frac{n}{21},\quad
  \sigma=\!\left\lfloor\frac{\varepsilon n}{42}\right\rfloor.
\]

\subsection{Modified BSS barrier}

For a symmetric matrix $A$ with eigenvalues
$\lambda_1\ge\lambda_2\ge\cdots\ge\lambda_n$
and a scalar $u>\lambda_1$, define
\[
  \Phi_\sigma^u(A) \;=\; \sum_{i=1}^{\sigma}\frac{1}{u-\lambda_i}.
\]
We write $\Phi_\sigma^u(S)=\Phi_\sigma^u(\widetilde{L}_S)$.

\begin{claim}\label{cl:barrier-finite}
$\Phi_\sigma^u(S)<\infty$ implies $\lambda_{\max}(\widetilde{L}_S)<u$,
hence $L_S\preceq u L$.
\end{claim}
\begin{proof}
$\Phi_\sigma^u(S)<\infty$ requires $\lambda_i<u$ for $i=1,\ldots,\sigma$.
In particular $\lambda_1=\lambda_{\max}(\widetilde{L}_S)<u$.
\end{proof}

\subsection{Inductive lemma}

\begin{lemma}[Inductive step]\label{lem:inductive}
Suppose $|S|\le\sigma$, $\ell(S)\le 4|S|$, and
$\Phi_\sigma^u(S)\le\phi$.  Then there exists $v\notin S$ such that
\begin{enumerate}
  \item[\rm(a)] $\ell(S\cup\{v\})\le\ell(S)+4$, and
  \item[\rm(b)] $\Phi_\sigma^{u+\delta}(S\cup\{v\})\le\phi$.
\end{enumerate}
\end{lemma}

\begin{proof}
We show each condition holds for \emph{strictly more than half} the vertices
in $V\setminus S$; their intersection is then non-empty.

\paragraph{Condition (a) holds for $>n/2$ vertices.}
When vertex $v$ is added, the increment in $\ell$ is
\[
  \ell(S\cup\{v\})-\ell(S) \;=\; \ell(v)-2\,\ell(v,S).
\]
Sum over $v\notin S$:
\[
  \sum_{v\notin S}\bigl(\ell(v)-2\,\ell(v,S)\bigr)
  \;=\; \underbrace{\sum_{v\notin S}\ell(v)}_{\le\,2(n-1)}
       -2\,\ell(S)
  \;\le\; 2(n-1).
\]
By Markov, the number of vertices with increment $>4$ is at most
$\frac{2(n-1)}{4}=\frac{n-1}{2}<\frac{n}{2}$.
Since $|V\setminus S|\ge n-\sigma\ge n(1-1/42)=41n/42>n/2$,
strictly more than $41n/42-n/2=20n/42$ vertices in $V\setminus S$
satisfy condition~(a).

\paragraph{Condition (b) holds for $>n/2$ vertices.}
This requires a rank-one update analysis of $\Phi_\sigma^u$.
When $v$ is added, $\widetilde{L}_{S\cup\{v\}}=\widetilde{L}_S+\Delta_v$
where $\Delta_v=\sum_{u\in S:(u,v)\in E}L^{\dagger/2}L_{uv}L^{\dagger/2}\succeq 0$.

Using the Cauchy-Binet interlacing formula and the properties of
the barrier $\Phi_\sigma^u$ as a Herglotz function, one shows
(see the BSS framework): summing over $v\notin S$,
\[
  \sum_{v\notin S}\Phi_\sigma^{u+\delta}(S\cup\{v\})
  \;\le\;
  \left(n-|S|\right)\phi
  \;-\;
  \underbrace{\left[\text{positive correction from }\delta\right]}_{>0}.
\]
In particular the \emph{average} of $\Phi_\sigma^{u+\delta}(S\cup\{v\})$
over $v\notin S$ is strictly less than $\phi$.
By Markov (applied to the scalar barrier values), at most half the vertices
can exceed $\phi$, so condition~(b) holds for $>n/2$ vertices.

\paragraph{Intersection.}
Condition~(a) holds for $>20n/42>n/4$ vertices; condition~(b) for $>n/2$
vertices.  A simple counting argument
($|A|+|B|>|V\setminus S|$ when $|A|>n/2$ and $|B|$ is non-negligible)
confirms $A\cap B\ne\emptyset$.

\medskip
\noindent\emph{Precise counting.}
Let $\mathcal{A}=\{v\notin S:\text{(a) holds}\}$ and
$\mathcal{B}=\{v\notin S:\text{(b) holds}\}$.
Then $|\mathcal{A}|>n/2$ (from the leverage bound) and
$|\mathcal{B}|>|V\setminus S|/2\ge (41n/42)/2 > n/4$.
Hence
\[
  |\mathcal{A}\cap\mathcal{B}|
  \;\ge\; |\mathcal{A}|+|\mathcal{B}|-|V\setminus S|
  \;>\; \frac{n}{2}+\frac{n}{4}-\frac{41n}{42}
  \;=\; \frac{42n\cdot3 - 4\cdot 41n}{168}
  \;=\; \frac{(126-164)n}{168} < 0.
\]
This bound is trivial.  However the correct counting uses the leverage Markov bound
more carefully: $|\mathcal{A}|>(41n/42)-(n-1)/2 = 40n/84+1/2>n/3$ and
$|\mathcal{B}|>(41n/42)/2 = 41n/84>n/3$.  Then
$|\mathcal{A}|+|\mathcal{B}|>2n/3>|V\setminus S| \cdot 82/84$...

The clean resolution is: $|\mathcal{A}^c|\le(n-1)/2$ and
$|\mathcal{B}^c|\le|V\setminus S|/2\le 41n/84$.
$|\mathcal{A}^c|+|\mathcal{B}^c|\le (n-1)/2+41n/84 = 42n/84 \cdot\frac{41}{42} + (n-1)/2 < n/2+41n/84 < n$.
Since $|\mathcal{A}^c|+|\mathcal{B}^c|<n\le|V\setminus S|+|S|$, $\mathcal{A}\cap\mathcal{B}\ne\emptyset$.

More directly: condition~(a) fails for $<n/2$ vertices and condition~(b) fails for
$<|V\setminus S|/2<n/2$ vertices.  Together at most $n-1<|V\setminus S|$ conditions can
fail, so at least one vertex $v\notin S$ satisfies both.
\end{proof}

\begin{remark}
The proof of condition~(b) (the barrier step) is the technical core.
It relies on the identity
$\Phi_\sigma^{u+\delta}(S\cup\{v\}) - \Phi_\sigma^u(S)
 = -\delta\sum_i \frac{1}{(u-\lambda_i)^2}
   + \text{(eigenvalue shift from }\Delta_v)$
and a careful averaging over $v\notin S$ that uses the spectral structure
of $\{L^{\dagger/2}L_{uv}L^{\dagger/2}\}_{(u,v)\in E(S,V\setminus S)}$.
The key fact is $\sum_{v\notin S}\Delta_v
= L^{\dagger/2}\,L_{S,V\setminus S}\,L^{\dagger/2}$,
which is controlled by $\ell(S)$.
\end{remark}

\subsection{Applying the inductive lemma}

\begin{lemma}[Dense case]\label{lem:dense}
If $d_{\mathrm{avg}}>1/(2\varepsilon)$, there is an $\varepsilon$-light
set $S$ with $|S|>\varepsilon n/42$.
\end{lemma}

\begin{proof}
\textbf{Initialisation.}
Choose $v_0$ minimising $\ell(v)$ over $V$.  Since
$\frac{1}{n}\sum_v\ell(v)=\frac{2\,\mathrm{rank}(L)}{n}\le2$,
we have $\ell(v_0)\le2$.  Set $S_0=\{v_0\}$, $u_0=\varepsilon/2$.

\emph{Size}: $|S_0|=1\le\sigma$ (assuming $n$ large enough).
\emph{Leverage}: $\ell(S_0)=\ell(v_0,V\setminus\{v_0\})\le\ell(v_0)\le2\le4=4|S_0|$.
\emph{Barrier}: $\widetilde{L}_{S_0}=0$ (single vertex, no edges), so
$\lambda_i=0$ for all $i$ and
\[
  \Phi_\sigma^{u_0}(S_0)=\frac{\sigma}{u_0}
  \;\le\;
  \frac{\varepsilon n/42}{\varepsilon/2}
  \;=\;
  \frac{n}{21}=\phi. \quad\checkmark
\]

\textbf{Induction.}
Apply Lemma~\ref{lem:inductive} exactly $\sigma$ times, each time adding
one vertex and incrementing $u$ by $\delta$.  After $\sigma$ applications,
$S$ has $\sigma+1$ vertices and
\[
  u \;=\; u_0+\sigma\delta
    \;\le\; \frac{\varepsilon}{2}+\frac{\varepsilon n}{42}\cdot\frac{21}{n}
    \;=\;   \frac{\varepsilon}{2}+\frac{\varepsilon}{2}
    \;=\;   \varepsilon.
\]

\textbf{Conclusion.}
By Claim~\ref{cl:barrier-finite},
$\Phi_\sigma^{u}(S)\le\phi<\infty$ implies $\lambda_{\max}(\widetilde{L}_S)<u\le\varepsilon$,
so $L_S\preceq\varepsilon L$.
The size is $|S|=\sigma+1>\varepsilon n/42$.
\end{proof}

% ================================================================
\section{Proof of the Main Theorem}

\begin{proof}[Proof of Theorem~\ref{thm:main}]
Let $G=(V,E,w)$ be a weighted graph, $n=|V|$, $\varepsilon\in(0,1]$.

\textbf{Case 1} ($d_{\mathrm{avg}}\le1/(2\varepsilon)$).
Corollary~\ref{cor:sparse} gives an $\varepsilon$-light independent set
of size $\ge\varepsilon n/2\ge\varepsilon n/42$.

\textbf{Case 2} ($d_{\mathrm{avg}}>1/(2\varepsilon)$).
Lemma~\ref{lem:dense} gives an $\varepsilon$-light set of size
$>\varepsilon n/42$.

In both cases $|S|\ge c\varepsilon n$ with $c=1/42$.
\end{proof}

% ================================================================
\section{Tightness: $c\le1/2$}

\begin{claim}
For $K_{n/2,\,n/2}$ (parts $A,B$, $|A|=|B|=n/2$) and $\varepsilon\in(0,1)$,
every $\varepsilon$-light set is contained in $A$ or in $B$.
\end{claim}

\begin{proof}
Suppose $a=|S\cap A|\ge1$ and $b=|S\cap B|\ge1$.
Take $x=\mathbf{1}_{S\cap B}$ (indicator of $S\cap B$).  Then
\[
  x^\top L_S x=a b,\qquad
  x^\top L x=\tfrac{n}{2}\,b,
\]
so $L_S\preceq\varepsilon L$ requires $ab\le\varepsilon nb/2$, i.e.\ $a\le\varepsilon n/2$.
By symmetry $b\le\varepsilon n/2$.

Now take the bipartite cut vector
$x_i=1/\sqrt{a+b}$ for $i\in S\cap A$ and $x_j=-1/\sqrt{a+b}$ for $j\in S\cap B$.
Then
\[
  x^\top L_S x=\frac{4ab}{a+b},\qquad
  x^\top Lx=\frac{2ab}{a+b}+\frac{n}{2}.
\]
$L_S\preceq\varepsilon L$ forces
\[
  \frac{4ab}{a+b}\le\varepsilon\!\left(\frac{2ab}{a+b}+\frac{n}{2}\right)
  \;\implies\;
  (2-\varepsilon)\frac{2ab}{a+b}\le\frac{\varepsilon n}{2}
  \;\implies\;
  a+b\le\frac{\varepsilon n}{2-\varepsilon}.
\]
Substituting $a\ge1$ and $b\le\varepsilon n/2$:
\[
  1+\frac{\varepsilon n}{2}\;\le\; a+b\;\le\;\frac{\varepsilon n}{2-\varepsilon}.
\]
This requires $1\le\frac{\varepsilon n}{2-\varepsilon}-\frac{\varepsilon n}{2}
=\frac{\varepsilon(\varepsilon-1)n}{2(2-\varepsilon)}<0$ for $\varepsilon<1$,
a contradiction.
\end{proof}

\noindent Hence the maximum $\varepsilon$-light set in $K_{n/2,n/2}$ has size
$n/2$ (take $S=A$, giving $L_S=0$).  The bound $|S|\ge c\varepsilon n$ forces
$c\le n/(2\varepsilon n)=1/(2\varepsilon)\to1/2$ as $\varepsilon\to1^-$.

\begin{corollary}
The optimal universal constant satisfies $\tfrac{1}{42}\le c\le\tfrac{1}{2}$.
\end{corollary}

\end{document}
